% !TeX root = ../dokumentation.tex

\addchap{\langanhang}

 
{\Large
\begin{enumerate}[label=\Alph*.]
	\item Code 
	\item Abbildungen

\end{enumerate}
}
\pagebreak
\*{A. Code}
\addcontentsline{toc}{section}{A. Code}
Der verwendete Code zum Auslesen der DHT11 Sensoren mit dem ESP32 und zum anzeigen der Messwerte auf einer Website.
\lstset{language=c++}
\begin{lstlisting}[caption={Code zur Auslesung des DHT11 Sensors mit dem ESP32}, label={lst:dht11_code}]
		#include <Arduino.h>
		#include "DHT.h"
		#include <WiFi.h>

		#define DHTPIN 14
		#define DHTTYPE DHT11

		const char* ssid = "iPhone VBackhaus";
		const char* password = "wtIhe357";

		DHT dht(DHTPIN, DHTTYPE);
		WiFiServer server(80);

		void setup() {
		Serial.begin(115200);
		
		WiFi.begin(ssid, password);
		Serial.println("Verbindet mit WLAN...");
		while (WiFi.status() != WL_CONNECTED) {
			delay(1000);
			Serial.println("Noch keine Verbindung...");
		}

		Serial.println("Verbunden");
		Serial.print("IP-Adresse: ");
		Serial.println(WiFi.localIP());

		dht.begin();
		server.begin();
		Serial.println("Server gestartet");
		}

		void loop() {
		WiFiClient client = server.available();
		if (!client) return;

		Serial.println("Neuer Client verbunden");
		String request = client.readStringUntil('\r');
		Serial.println(request);
		client.flush();

		//pruefen ob die Anfrage nach "/data" geht
		if (request.indexOf("GET /data") >= 0) {
			float h = dht.readHumidity();
			float t = dht.readTemperature();

			client.println("HTTP/1.1 200 OK");
			client.println("Content-Type: application/json");
			client.println("Connection: close");
			client.println();

			if (isnan(h) || isnan(t)) {
			client.println("{\"error\": \"Sensorfehler\"}");
			} else {
			client.printf("{\"temperature\": %.2f, \"humidity\": %.2f}", t, h);
			}
			return;
		}

		//HTML datei mit Script
		String html = R"rawliteral(
		<!DOCTYPE html>
		<html>
		<head>
			<meta charset="UTF-8">
			<title>DHT11 Sensor</title>
		</head>
		<body>
			<h1>ESP32 DHT11 Sensorwerte</h1>
			<p>Temperatur: <span id="temp">--</span> C</p>
			<p>Luftfeuchtigkeit: <span id="hum">--</span> %</p>

			<script>
			// Diese Funktion holt neue Daten vom ESP32
			async function ladeDaten() {
				const antwort = await fetch('/data');      // Anfrage an /data schicken
				const daten = await antwort.json();        // Antwort als JSON lesen
				document.getElementById('temp').textContent = daten.temperature;
				document.getElementById('temp').textContent = daten.temperature;
			}

			ladeDaten();                 // Beim Laden der Seite einmal aufrufen
			setInterval(ladeDaten, 3000); // Alle 3 Sekunden automatisch neu laden
			</script>
		</body>
		</html>
		)rawliteral";

		client.println("HTTP/1.1 200 OK");
		client.println("Content-type:text/html");
		client.println("Connection: close");
		client.println();
		client.println(html);
		client.println();
		}
\end{lstlisting}
% Noch machen: Code einfügen (was zu hölle läuft hier falsch mit listings???)

\section*{B. Abbildungen}

\begin{figure}[H]
	 \centering 
	 \resizebox{\textwidth}{!}
	 { \includegraphics{circuit_image.png} } 
	 \vspace{10pt} 
	 \caption{Anhang: Schaltplan DHT11 (eigene Abbildung)}\label{anhang: Schaltplan} 
\end{figure}

\begin{figure}[H]
	 \centering 
	 \resizebox{\textwidth}{!}
	 { \includegraphics{Luftfeuchtigkeitsmessreihe.png} } 
	 \vspace{10pt} 
	 \caption{Anhang: Luftfeuchtigkeitsmessreihe}\label{anhang: Luftfeuchtigkeitsmessreihe} 
\end{figure}


\begin{figure}[H]
	 \centering 
	 \resizebox{\textwidth}{!}
	 { \includegraphics{Luftfeuchtigkeitsmessreihe.png} } 
	 \vspace{10pt} 
	 \caption{Anhang:Temperaturmessreihe}\label{anhang: Temperaturmessreihe} 
\end{figure}

\begin{figure}[H]
	 \centering 
	 \resizebox{\textwidth}{!}
	 { \includegraphics{ausschnitt_Taupunkttabelle.png} } 
	 \vspace{10pt} 
	 \caption{Anhang: Auschnitt aus der Taupunkttabelle \protect\citeM{d}{tauKor}{1}{1}}\label{anhang:Auschnitt_Taupunkttabelle} 
\end{figure}



\pagebreak
%\includepdf[pages=-,scale=.9,pagecommand={}]{Aufgabenstellung.pdf} % PDF um 10% verkleinert einbinden --> Kopf- und Fußzeile  werden so korrekt dargestellt. Die Option `pages' ermöglicht es, eine bestimmte Sequenz von Seiten (z.B. 2-10 oder `-' für alle Seiten) auszuwählen.
\pagebreak



